%% kapa: exam=1 round=21 year=2010 period=1 form=A subject=civil_law subject_ko=민법 seq=15 number=15
%% source: Q-Net 감정평가사 기출문제 (cst00309) artlSeq=5024181 file=extracted/1008/5024181/2010년 제21회 감정평가사 1차시험 기출문제/1교시.hwp
%% conversion: kordoc 4.24.51 -> tools/gen1.py (auto); review=proofread-sonnet; flags=none
\begin{question}{15}
甲의 대리인 乙이 대리행위로서 丙으로부터 토지를 매수하는 계약을 체결하였다. 다음 설명으로 옳지 않은 것은?
\begin{choices}
  \item 대리행위의 효력이 의사의 흠결로 영향을 받을 경우, 그 사실의 유무는 甲을 표준으로 하여 결정한다.
  \item 乙이 甲을 위한 것임을 표시하지 않았다면, 그 계약은 乙을 위한 것으로 보는 것이 원칙이다.
  \item 乙이 미성년자이더라도 丙과 체결한 계약의 효력은 甲에게 생긴다.
  \item 乙이 계약을 체결하기 전에 이미 파산한 경우, 丙이 선의ㆍ무과실인 때에는 甲에게 그 계약의 효력을 주장할 수 있다.
  \item 만약 乙에게 대리권이 없는 경우, 丙이 계약을 철회하기 전에 甲이 추인하면 매매계약의 효력은 甲에게 생긴다.
\end{choices}
\end{question}
\answer{1}
